% Opciones empleadas:
%
%	a4paper -> indica el tama�o del papel, en este caso A4.
%	11pt -> tama�o de la fuente 11 puntos.
%	oneside -> s�lo escribimos en una cara del folio.
% 
% Otras opciones interesantes:
%
%	twoside -> escribimos a doble cara.
%	openbib -> para que las referencias bibliogr�ficas tengan un salto de l�nea entre cada campo de la referencia.
%
\documentclass[a4paper,11pt,twoside]{book}

% codificaci�n latin1 y s�mbolos del idioma espa�ol (�, acentos, ...)
\usepackage[english]{babel}
\usepackage[latin1]{inputenc}

% puede que queramos usar el s�mbolo del euro.
\usepackage{eurosym}

% El paquete fancybox nos permite crear cajas de diferentes estilos con facilidad.
% http://www.ctan.org/get/macros/latex/contrib/fancybox/fancybox.pdf
% http://www.mackichan.com/index.html?techtalk/487.htm~mainFrame
\usepackage{fancybox}
\usepackage{multicol}
\usepackage{amsmath}

% Para incluir subfiguras.
\usepackage{subfigure}

% Para incluir gr�ficos en JPG => compilar con pdflatex.
% \usepackage[pdftex]{graphicx}
\usepackage[pdftex]{graphicx}   % --> LaTex =>PDF con epstodf
\usepackage{epstopdf}

% Para incluir gr�ficos EPS => compilar con latex.
%\usepackage[dvips]{graphicx}

% Para escribir en color...
%
% ... cuando compilamos con el comando ``latex''
%\usepackage[dvips,usenames]{color}
% ... uando compilamos con el comando ``pdflatex''
\usepackage[pdftex,usenames,dvipsnames]{color}

% Espaciado y ajuste de m�rgenes
\usepackage{setspace}
\onehalfspacing
% \doublespacing
\setlength{\textwidth}{15cm}
\setlength{\textheight}{22cm}
\setlength{\oddsidemargin}{1cm}
\setlength{\evensidemargin}{1cm}
%\usepackage[hmargin={4cm,2.5cm},tmargin=4cm,bmargin=4cm]{geometry}
%Code snippets.
\usepackage{listings}
%Landscape figures.
\usepackage{lscape}

% Paquete fancyhdr -> Para modificar la cabecera y pie de p�ginas.
% http://tug.ctan.org/tex-archive/macros/latex/contrib/fancyhdr/
\usepackage{fancyhdr}
\pagestyle{fancy}
\fancyhf{}
%\fancyhf[HR]{\thepage}
%\fancyhf[HL]{\nouppercase\rightmark}
\fancyhead[LO]{\leftmark}
\fancyhead[RE]{\rightmark}
\fancyhead[RO,LE]{\thepage}
\renewcommand{\chaptermark}[1]{\markboth{\textbf{\thechapter. #1}}{}}
\renewcommand{\sectionmark}[1]{\markright{\textbf{\thesection. #1}}}
\renewcommand{\headrulewidth}{0.6pt}
\setlength{\headheight}{2.5\headheight}

% Package booktabs -> Para mejorar el aspecto de las tablas o cuadros.
% http://www.ctan.org/tex-archive/macros/latex/contrib/booktabs/
\usepackage{booktabs}

% Package rotating -> Para poder girar las tablas y dibujarlas a lo largo
% del folio en vez de a lo ancho.
\usepackage{rotating}

% Packages multicol y multirow, para manejar tablas de filas y columnas m�ltiples.
\usepackage{multicol}
\usepackage{multirow}
%To center the captions.
\usepackage[center]{caption}

%Links in biblio.
\usepackage{url}

\usepackage[pagebackref=true, hypertexnames=true, plainpages=false]{hyperref}

%Bold colors
\usepackage{color,soul}

%Math caligraphy
\DeclareMathAlphabet{\mathpzc}{OT1}{pzc}{m}{it}

%\usepackage[T1]{fontenc}

% Personalizamos la separaci�n entre p�rrafos...
\parskip=6pt

% Personalizamos el identado en la primera l�nea del nuevo p�rrafo...
\parindent=0pt

% Establecemos el n�mero m�ximo de niveles de profundidad en las secciones.
\setcounter{secnumdepth}{3}

%C++ symbol
\def\CC{{C\nolinebreak[4]\hspace{-.05em}\raisebox{.4ex}{\tiny\bf ++}}}

% T�tulo
\title{Design and implementation of a HPC solution for stereo matching}
% Autor
\author{Luis Omar Alvarez Mures}
% Fecha
\date{\today}


\begin{document}

	% \maketitle sirve para generar autom�tica una portada predefinida, pero para un proyecto fin de carrera
	%	de FIC no sirvir�a porque no cumple las normas de presentaci�n. Podemos hacer dos cosas:
	% 1. Usarla e ignorar las normas (y asumir las consecuencias que pueda tener)
	% 2. Hacernos una portada en LaTeX que cumpla las normas (menos arriesgado)
	%
	
	%Modify list bullets so it does not use dash for second level
	\renewcommand{\labelitemi}{$\bullet$}
	\renewcommand{\labelitemii}{$\circ$}
	
	%Change page numbering so backref works
	\pagenumbering{roman}	
	
        %
% Portada.
%

% Nota: Ser�a m�s c�modo emplear el comando \maketitle que genera una portada de forma autom�tica, pero 
% no incluye toda la informaci�n que es necesario incluir en la memoria de un proyecto de fin de carrera
% de la Facultad de Inform�tica de A Coru�a.
%

\begin{titlepage}

  \begin{center}
    \includegraphics[width=7cm]{figures/anagramaUDC.pdf}\\[0.75cm]
    {\textsc{Faculty of Computer Science}} \\
    {\large \textsc{Departament of Electronics and Systems}} \\[1cm]
    {\Large \textsc{Final Year Project}} \\[0.25cm]
    {\Large \textsc{Master in High Performance Computing}} \\[2cm]
    {\LARGE \textsl{\textbf{Design and implementation of a HPC solution for stereo matching}}} 
    \vfill
    \begin{flushright}
      \begin{tabular}{ll}
        \textbf{Student:}    & Luis Omar �lvarez Mures\\
        \textbf{Advisors:} & Emilio J. Padr�n Gonz�lez\\
        & Juan Ram�n Rabu�al Dopico\\
        & \\
        \multicolumn{2}{r}{\small \emph{A Coru�a, \today{}.}} \\
      \end{tabular}
    \end{flushright}
  \end{center}

\end{titlepage}

        \thispagestyle{empty}
        \cleardoublepage
        %\newpage
	%\thispagestyle{empty}
	%\mbox{}

	% FRONTMATTER: TOC, LOF, LOT y descripci�n/organizaci�n de la memoria.
        \frontmatter
	
	% Los proyectos de fin de carrera de FIC han de ir acompa�ados de una serie de documentos adicionales, algunos
	% 	de ellos obligatorios (certificado, resumen, lista de palabras clave) y otros opcionales (dedicatoria
	%	y agradecimientos).
	%
	
        \thispagestyle{empty}     % No number page, headings...
	\thispagestyle{empty}
\section*{Informaci�n xeral}
\vfill
\begin{center}
  \begin{tabular}{p{4.5cm}p{9cm}}
    \textbf{\emph{T�tulo do proxecto:}} & \textbf{``Dese�o e implementaci�n dunha soluci�n HPC para Stereo Matching''} \\[0.5cm]
    \emph{Clase de proxecto:} & Proxecto de desenvolvemento en investigaci�n \\[0.5cm]
    \emph{Nome do estudante:} & Luis Omar �lvarez Mures \\[0.5cm]
    \emph{Nome dos directores:} & Emilio Jos� Padr�n Gonz�lez \\
    & Juan Ram�n Rabu�al Dopico \\[0.5cm]
    \emph{Membros do tribunal:} & \\[0.5cm]
    & \\
    & \\
    & \\
    & \\
    & \\
    & \\
    & \\
    & \\
    & \\
    & \\
    \emph{Data de lectura:} & \\[0.5cm]
    \emph{Cualificaci�n:} & \\
  \end{tabular}
\end{center}
\vfill

	
	\thispagestyle{empty}
        \cleardoublepage
	
	\thispagestyle{empty}     % No number page, headings...
	%
% Dedicatoria
%
\section*{}

\begin{flushright}
	{\it Dedicated to Rebeca and my parents.}
\end{flushright}

        
        \thispagestyle{empty}
        \cleardoublepage
        
        %\newpage
	%\thispagestyle{empty}
	%\mbox{}
	
	\thispagestyle{empty}
        %
% Agradecimientos
%

\section*{Acknowledgements\footnote{In alphabetical order}}

My sincerest thank you for the help provided:

To Mr. Emilio Jos� Padr�n Gonz�lez.

To Mr. Juan Ram�n Rabu�al Dopico.




        
        \thispagestyle{empty}
        \cleardoublepage
        %\newpage
	%\thispagestyle{empty}
	%\mbox{}
	
	\thispagestyle{empty}
        %
% Resumen del proyecto de fin de carrera
%

\section*{Resumo:}

O obxectivo deste proxecto � o dese�o e implementaci�n dunha
ferramenta multiplataforma para a visualizaci�n de nubes de puntos que
permita a xesti�n e manipulaci�n de grandes volumes de datos. Este
tipo de \emph{datasets} poden ser obtidos a partir dun esc�ner LIDAR
ou de alternativas de baixo custo como KINECT.

Este tipo de sistemas capturan un gran volume de informaci�n, xerando
nubes de puntos de gran tama�o con miles de mill�ns de informaci�n
xeom�trica, color, normais, reflectividade, etc. A consecuencia � un
conxunto de informaci�n dif�cil de manexar de xeito interactivo, ao
non caber completo na memoria RAM dun ordenador e requerir o traballo
desde memoria secundaria (cl�sicamente un disco duro). Isto complica
a creaci�n de aplicaci�ns en tempo real para a xesti�n e manexo deste
conxunto masivo de datos.

A ferramenta software desenvolvida neste proxecto ofr�celle ao usuario
final todas as funcionalidades necesarias para o traballo con este
tipo de datos, desde un completo visualizador 3D para render de
puntos, que incl�e caracter�sticas de visualizaci�n avanzada de
OpenGL, soporte para o traballo con varias nubes de puntos
simultaneamente, m�ltiples c�maras diferentes coas que traballar,
caracter�sticas multiresoluci�n, posibilidade de aplicar
transformaci�ns e distintos tipos de pre-proceso sobre os datos, medir
distancias, detecci�n de obxectos na nube de puntos, exportaci�n a CAD
dos obxectos seleccionados, etc

Para a implementaci�n de todas esas funcionalidades partiuse da
biblioteca software PCM, a�nda en desenvolvemento e que xurde tam�n
na Unidersidade da Coru�a. Este biblioteca proporciona unha serie de
ferramentas de baixo nivel para a xesti�n de nubes de puntos masivas
en hardware de uso com�n. Neste proxecto non soamente se fai uso das
caracter�sticas existentes de PCM, sen�n que, ademais da propia
ferramenta de visualizaci�n desenvolvida (bautizada como ToVIEW), tam�n se extendeu
ampliamente a biblioteca PCM para poder ofrecer todas as funcionalidades
necesarias no proxecto, ademais de varias melloras de rendemento na mesma.

No proxecto abordouse ademais o aproveitamento das posibilidades para
o procesamento paralelo presentes no hardware no que se executa PCM,
tanto en canto � CPU como �s capacidades GPGPU da tarxeta gr�fica.

Por �ltimo, o paquete software resultante foi dese�ado para ser
facilmente integrable nos fluxos de traballo habituais nos sectores da
enxe�er�a que m�is habitualmente empregan escaneos 3D en forma de
nubes de puntos, como pode ser, por exemplo, a topograf�a.
Durante a realizaci�n do proxecto traballouse con profesionais da
topograf�a habituados a incorporar o uso da tecnoloc�a LIDAR no seu
traballo, o que facilitou a creaci�n dunha ferramenta �til e vers�til
para o usuario final.
No aspecto da f�cil integraci�n en entornos profesionais tam�n axuda o
feito de tratarse dun desenvolvemento multiplataforma e \emph{open source}.

        
        \thispagestyle{empty}
        \cleardoublepage
	
	\thispagestyle{empty}
        %
% Resumen del proyecto de fin de carrera
%

\section*{Abstract:}

The objective of this project is designing and implementing a multi-platform point-based rendering visualizer that allows the management and interactive manipulation of big 3D point clouds. These clouds can be obtained using a LIDAR scanner or cheaper alternatives like KINECT. 

These systems can nowadays generate data-sets with huge amounts of point data, position, color, normals, reflectivity, etc. As a consequence, these point clouds will not fit on system RAM and will have to reside in an HDD. Due to this fact, real-time rendering of massive point clouds is a complex task.

The resulting software tool developed in this project, will offer the end user the functionality necessary for working with these types of datasets. Including a complete 3D visualizer, with advanced OpenGL point-rendering capabilities, multiple point cloud support, multiple cameras, multi-resolution capabilities, point cloud transformation and preprocessing, distance measurements, object segmentation, CAD exportation of segmented objects, etc. 

To implement all of these features, we will utilize PCM. PCM is a project in development from the University of A Coru�a, that provides a series of low level tools for working with point clouds on commodity hardware. This project will not only extend PCM with the aforementioned functionality and a high level software tool for the end user, but will improve the existing code and some performance aspects of PCM.

Furthermore, the provided tool will take advantage of all the parallel computing capabilities of the target platforms. Using multiple cores on CPUs and GPGPU on the graphics cards. The multi-resolution framework will also help when trying to achieve maximum performance when trying to apply any computation.

In conclusion, the resulting software package should be easily integrated and improve substantially the existing engineering workflows, thanks to the fact that the tool will be multi-platform and open source. Since it was also developed with the advice of some surveyors that are working with laser scanners everyday, there are high hopes of truly having created a versatile and useful tool.





        
        \thispagestyle{empty}
        \cleardoublepage
        
        \thispagestyle{empty}     % No number page, headings...
	%
% Palabras clave
%

\section*{Lista de palabras chave:}

Motor, render, nubes de puntos, k-d tree, multi-resoluci�n, tiempo real, PCM, ToView, open source, multi-plataforma.

\section*{Keywords:}

Engine, render, point clouds, k-d tree, multi-resolution, real-time, PCM, ToView, open source, multi-platform.



	
	\thispagestyle{empty}
        \cleardoublepage
	
        \thispagestyle{empty}     % No number page, headings...
        %
% Software y Hardware
%

\section*{Hardware e software empregado:}

\begin{itemize}
\item Intel Core i7s, NVIDIA GTXs \& GTs.
\item C++, Visual Studio, Git, OpenCV, OpenMP, OpenCL, CUDA, Boost, GCC.
\end{itemize}



        
        \thispagestyle{empty}
        \cleardoublepage
        
       % \thispagestyle{empty}     % No number page, headings...
        
        %\newpage
	%\thispagestyle{empty}
	%\mbox{}

        \tableofcontents
        \listoffigures
        \listoftables
        
        %\newpage
	%\thispagestyle{empty}
	%\mbox{}

	% MAINMATTER: El contenido, cap�tulo a cap�tulo, de la memoria del tfg.
        \mainmatter
	
	\pagenumbering{arabic}	
	
	%
% Frontmatter - Introducci�n. Los miembros del tribunal que juzgan los PFC's tienen muchas m�s memorias que leer, por lo que
%	agradecer�n cualquier detalle que permita facilitarles la vida. En este sentido, realizar una peque�a introducci�n,
%	comentar la organizaci�n y estructura de la memoria y resumir brevemente cada cap�tulo puede ser una buena pr�ctica
%	que permita al lector centrarse f�cilmente en la parte que m�s le interesa.
%

\chapter[Introduction]{
	Introduction
}

Computer Graphics is the discipline that deals with creating, managing, analyzing and visualizing graphics using a computer, and it is closely linked to other disciplines such as Image Processing, Computer Vision or Machine Learning, among others. Advances in computer graphics have had an impact on several media types such as, for example, the movies and video games industry, civil engineering, medicine, etc.

In computer graphics, \emph{rendering} is the process of generating a 2D image to be shown in a display. This process takes usually a virtual scene of 3D objects as input, containing data like geometry, lighting, cameras, etc.
Typically, polygonal modeling has been the approach followed for modeling objects in the scene, approximating their surfaces using polygons. Point-based graphics focuses on points as the fundamental representation of the surfaces instead of polygons (see \figurename~\ref{poly_comp}).

\begin{figure}
	\centering
	\includegraphics[scale=0.6]{figures/pol_comp.pdf}
	\caption[Modeling types comparison]{
		From left to right; mathematical representation of a plane, representation using two polygons (triangles) and a point-based representation.
	}
	\label{poly_comp}
\end{figure}

Nowadays, acquisition technologies such as LIDAR (Laser Imaging
Detection and Ranging) allow us to obtain high-precision
georeferenced 3D scans of the real world in an easy and fast way.
These systems use ultraviolet, visible or near infrared light to image objects. It can target a wide range of materials and even single molecules. A narrow laser beam can map features with very high precision.
The evolution of this kind of technologies is having a tremendous
impact in multiple fields in the engineering, and bringing back the
camp-work to the office has become a standard practice.
This presents clear advantages compared to the classical work-flow in
disciplines like topography, that usually need to come back on the
ground each time a new measurement is required. A couple of images
obtained from a LIDAR scanner are shown in
\figurename~\ref{lidarexamples}.

\begin{figure}
  \centering
  \includegraphics[width=0.475\textwidth]{figures/campillo1.png} \quad
  \includegraphics[width=0.475\textwidth]{figures/parcela.png}
  \caption{Examples of LIDAR scanning.}\label{lidarexamples}
\end{figure}

These scanning devices generate data-sets with an enormous number of
samples, with 3D coordinates and additional data such as color,
reflectivity and so on. These point clouds, as they are usually
named, are formed by hundreds or thousands of points, resulting in
Giga or even tera-bytes of information.
Using point based representations, the polygonal approximation step
can be skipped and the massive amount of captured data is directly
used. However, managing these huge data-sets, either for applying some kind of data
analysis or for visualization, is a computational challenge and an
open and state-of-the-art problem. This project deals with this question.

This chapter introduces the motivation behind the project and the
objectives we try to achieve. It shows the structure of this document
as well.

\section[Motivation and context]{
	Project motivation and context
}

PCM is a software library developed in the University of A Coru�a for the management of huge point datasets obtained by LIDAR in commodity hardware systems. It was started as a MSc Final Project~\cite{Jaspe12} (system structure and basic capabilities) in the Faculty of Computer Science and then continued in a BSc Final Project~\cite{Jaspe13} (improved render and other minor additions).

PCM is based on a multi-resolution structure and includes basic memory hierarchy management, making it possible to transfer the point cloud data between secondary memory (HDD) and video memory in the graphics card. Originally, PCM includes a very basic visualizer, quite simple and limited but that has the potential to become the seed for a more advanced point cloud software tool that will allow the user to access all of these features in real-time.

Although PCM was designed to deal with arbitrary size datasets, it is far from being a mature software and has some problems managing point clouds in the order of billion points. Because of the massive nature of these datasets, achieving real-time manipulation of them is a complex task. It will require the improvement of PCM, with a complete profiling and performance analysis and a careful refactoring. Having a system that is able to manage such a huge amount of data, will let the user take advantage of the full potential of cutting-edge 3D capture devices.
Of course, the library should exploit the full potential of both modern CPUs and GPUs, automatically using parallel processing to reduce the computation time as much as possible.
Furthermore, although PCM sketches a framework for applying data
analysis on the point clouds, this part was barely crafted at
all. Other constraints of the original proposal are also addressed in
this work, such as adding read/write operations to allow the
modification of the datasets, the possibility of working with multiple
point clouds, etc.

In this project, a complete user-oriented software suite for the visualization, management, analysis and manipulation of massive 3D point clouds will be developed from the original PCM proposal. ToView is the name chosen for this software package. It is a multi-platform (GNU/Linux, Windows and MacOS) free software project and makes use of open source software exclusively. For these reasons the user interface is based on the cross-platform application framework QT, that has native implementations and tools for all the aforementioned platforms. The resulting QT application will let the user access all of the functionality including an advanced and complete 3D visualizer.

\section[Objectives]{
	Project objectives
}

The aim of this project is the design and implementation of a point-based multi-platform software tool for the interactive visualization, management, analysis and manipulation of massive 3D point clouds. 

The application will offer the user the necessary tools to work effectively with this type of datasets, including a 3D visualizer with advanced point-rendering techniques using OpenGL, that will allow real time interaction with one or multiple point clouds.

From the interface, the user will be able to:

\begin{itemize}
\item Select different visualization modes for point-based models: multiple cameras, simple and advanced visualization, multi-resolution, etc.
\item Combine different point clouds with tools that enable rotation, translation and scaling of the datasets.
\item Select points in the clouds for working with them.
\item Apply operations to a part of the cloud or the complete dataset. From simple operations, such as distance measurement and calculation, to plane or other geometric primitives detection.
\item Export the work done to a standard CAD format, so the tool can be integrated in other work-flows.
\item Call the pre- and post-processing tools included in the suite: multi-resolution structure computation, filters, etc.
\end{itemize}

The developed tool will take advantage of the parallel computing possibilities of the platforms. Both using the GPGPU capabilities of the graphics card with OpenCL and the multiple cores of the CPU. 

PCM will be used as the foundation for this work. PCM is a project in development that offers several low level tools for working with point clouds of arbitrary size in commodity hardware, from which a multi-resolution structure with different levels of cache is highlighted. This final degree project will not only extend PCM with the mentioned features, offering a high level software suite for the final user, but will also improve the existing source code, paying special attention to performance aspects.

A series of benchmarks will also be implemented that will allow us to obtain automatically and autonomously performance information about PCM and the visualization tool, with graphs and reports. This tool will make finding bottlenecks and performance issues easier.

\section[Structure]{
Document structure
}

The rest of the document is organized following this structure:

\paragraph*{Chapter 2.}
Planning and methodology

After this introduction, we start by describing the planning and
development methodology used in the project. We chose an agile
software development approach.

\paragraph*{Chapter 3.}
Computer Graphics basics

This chapter briefs the very basic concepts about computer graphics needed to understand the rest of the document: How a virtual camera works and the different types of cameras, what 3D transformations are and how a point is represented in 3D.
 
\paragraph*{Chapter 4.}
Structure of a real-time visualizer for point-based datasets

This chapter describes the structure of the visualizer we propose in this work and presents the class diagram of the software package. Both ToView (the visualization tool) and PCM (the multi-resolution library behind the GUI) are included.

\paragraph*{Chapter 5.}
Point Cloud Manager

This chapter focussed in describing PCM ant detailing all the work
invested in the improvement of the original PCM library and tools:
code refactoring, bug fixes and new functionalities. It includes the
experimental results of the tests included in the benchmarks we have
designed to measure the library performance.

\paragraph*{Chapter 6.}
ToView

This chapter depicts the visualizer designed and implemented from
scratch in this project to manage and visualize massive point
clouds. All the included functionalities are described and
experimental results are also included.

\paragraph*{Chapter 7.}
Advanced GPU point rendering

This chapter is devoted to the most advanced visualization methods
included in ToView, a bunch of some of the most state-of-the-art
point-rendering techniques, and how they are implemented with OpenGL
in the GPU using GLSL (OpenGL Shading Language).

\paragraph*{Chapter 8.}
Object detection: RANSAC

This chapter deals with the detection of different shapes and volumes
in ToView, describing the different approaches that were studied to implement
this functionality in the project. Specifically, an extensive
explanation of the RANSAC algorithms used in this project is provided.

\paragraph*{Chapter 9.}
Preprocessing and filtering

Both PCM and ToView need several pre-processing tasks to take advantage
of some of the most advances features. All this pre-processing and the
different filtering techniques applied to the datasets are explored in
this chapter.

\paragraph*{Chapter 10.}
Conclusions and future lines of work

Finally, this chapter presents the main conclusions obtained with this
project and outlines possible extensions to both PCM and ToView.

	%\include{tfg_main_020}
	%\include{tfg_main_030}

	% INCLUIMOS LOS AP�NDICES...
        %\appendix
	%\include{tfg_appendix_010}
%         \include{tfg_appendix_020}

	%Primero hay que compilar el bibtex para que salga.
	% INCLUIMOS LA BIBLIOGRAF�A...
	\nocite{*}	% Se usa para indicar en la bibliograf�a las referencias no citadas.
	\bibliography{pfm_biblio}
	\bibliographystyle{alpha}

\end{document}

